\section{实验结果与分析}

本章呈现所有定量实验结果与定性案例分析。下述数据均来自测试集（共计 1997 个 trial）的完整评估，A2 模型的结果为三个随机种子（42、123、2025）的平均值。

\subsection{实验设置}

本实验所采用的软硬件环境汇总于表~\ref{tab:exp-setup}。

\begin{table}[htbp]
  \centering
  \caption{实验软硬件配置}
  \label{tab:exp-setup}
  \begin{tabular}{l l}
    \toprule
    配置项 & 说明 \\
    \midrule
    GPU & NVIDIA RTX 6000 Ada, 48GB 显存 \\
    精度 & bfloat16 混合精度 \\
    CLIP backbone & ViT-B/32（pooled 512 维, visual tokens $50\times512$）\\
    EEG encoder & A2 时间-频谱-空间编码器, 512 维输出 \\
    训练规模 & train 7970, val 1998（全量训练）\\
    测试规模 & test 1997（全量测试）\\
    退化条件 & clean, lowres16, strong\_blur, strong\_noise, occlusion50, mixed \\
    EEG 模式 & vision\_only, real\_eeg, shuffled\_eeg, random\_eeg, eeg\_only \\
    \bottomrule
  \end{tabular}
\end{table}

\subsection{A2 语义识别结果}

表~\ref{tab:a2-main} 列出了 A2 模型在六种图像退化条件下的多随机种子语义分类结果，所有数值均为三个种子运行得到的平均 top-1 准确率。

\begin{table}[htbp]
  \centering
  \caption{A2 语义融合多种子语义预测结果（三种子均值）}
  \label{tab:a2-main}
  \begin{tabular}{l c c c c c c c}
    \toprule
    退化条件 & Vision-only & Real EEG & Shuffled EEG & Random EEG & Real-Vision & Real-Shuffled & Real-Random \\
    \midrule
    clean & 0.950 & 0.975 & 0.403 & 0.681 & +0.025 & +0.572 & +0.294 \\
    lowres16 & 0.258 & 0.523 & 0.070 & 0.062 & +0.265 & +0.453 & +0.461 \\
    mixed & 0.066 & 0.392 & 0.046 & 0.041 & +0.326 & +0.346 & +0.351 \\
    occlusion50 & 0.811 & 0.880 & 0.237 & 0.378 & +0.069 & +0.643 & +0.502 \\
    strong\_blur & 0.759 & 0.837 & 0.187 & 0.273 & +0.078 & +0.650 & +0.564 \\
    strong\_noise & 0.644 & 0.782 & 0.178 & 0.266 & +0.139 & +0.605 & +0.516 \\
    \bottomrule
  \end{tabular}
\end{table}

观察表~\ref{tab:a2-main} 可得出以下发现。第一，real EEG 在所有退化情形下的准确率均高于 vision-only，增益幅度从 clean 条件下的 2.5 个百分点到 mixed 条件下的 32.6 个百分点不等，退化愈严重则提升愈明显。第二，real EEG 在所有退化条件下均显著优于 shuffled EEG（打乱脑电）和 random EEG。在 lowres16、mixed、strong\_blur、strong\_noise 等强退化条件下，shuffled EEG 和 random EEG 的性能远不及 real EEG；而在 clean 条件下，对照组中的视觉分支依然能够提供较好的分类准确率，因此其与 real EEG 的差距相对较小。这说明对照组的表现并非纯粹的随机猜测——视觉输入仍然发挥作用——只是 EEG 与图像之间的配对关系被破坏后，其分类能力回落至接近纯视觉基线的水平。控制实验表明，性能增益并非来源于额外的输入通道或更大的模型容量，而是依赖于 EEG 与图像之间的真实对应关系。第三，在全部 3 个种子乘以 6 种条件共 18 组评估中，real EEG 的表现在每一种设定下均领先于所有对照条件。

\placeholderfigure{A2 不同退化条件下 Top-1 accuracy 柱状图}{a2-bar-chart-placeholder}{5.0cm}

\subsection{不同语义融合方法对比}

除了 A2 之外，本设计还考察了若干其他融合策略：(1) P2 raw+spectrogram 双路融合——仅组合时间波形与频谱两路特征，不含空间分支；(2) P2A2——在 P2 编码器基础上增设 A2 风格的三支路后进行融合；(3) residual adapter——在视觉嵌入上附加 EEG 残差适配器，不使用门控机制；(4) prototype bias——在文本原型相似度计算中引入 EEG 预测的偏移量。整体而言，A2 在大多数退化条件下取得了最大的 real-vision 差距，且 real EEG 相对于 shuffled EEG 和 random EEG 的优势最为一致。三支路编码器配合门控残差融合的设计思路已在第 5 章中阐述，实验数据初步印证了该方案的有效性。

\subsection{Token 级 VTF 结果分析}

VTF 在 token 级别上实现了 real EEG 超越 vision-only 及 EEG 对照组的分类效果，说明 token 级的 EEG 融合方案具有可行性。然而，VTF 的绝对分类准确率明显低于 A2 的 pooled 融合方式（例如在 lowres16 条件下，VTF 的 real EEG top-1 准确率约为 0.18，而 A2 达到 0.52）。这一差异可能源于两方面原因。其一，A2 直接在 CLIP 语义空间中对类别预测进行优化，任务目标更为直接；而 token 级融合在小样本场景下训练难度较大，分类稳定性不及 embedding 层面的 A2。其二，VTF 仅采用单层 cross-attention 机制，而 A2 使用了更深层的门控融合网络，信息交互更加充分。VTF 的核心价值在于其输出的增强视觉 token 可以无缝接入后续的生成式模型，是为满足生成任务需求而设计的关键中间结构。

\subsection{生成式视觉语言模型结果}

各个生成方案的定量评估结果汇总于表~\ref{tab:evlm-results}。其中 Class-hit 一列表示在 strong\_blur 和 strong\_noise 两个强退化条件下 caption 类别命中率的平均值。

\begin{table}[htbp]
  \centering
  \caption{生成式视觉语言模型结果对比}
  \label{tab:evlm-results}
  \begin{tabular}{l c c c c c c}
    \toprule
    模型 & 有效率 & 无效率 & Class-hit & Real-Vision & Real-Shuffled & Real-Random \\
    \midrule
    GRU baseline & 68.1\% & 31.9\% & 12.4\% & +2.2\% & +5.2\% & +4.1\% \\
    方案一 Qwen2.5 + LoRA & 32.4\% & 67.6\% & 18.1\% & +2.7\% & +5.3\% & +5.3\% \\
    方案五 full adapter & 91.1\% & 8.9\% & 29.4\% & +5.0\% & +5.5\% & +3.0\% \\
    方案五 LoRA $r=8$ & 91.9\% & 8.1\% & 38.3\% & +9.1\% & +10.3\% & +7.8\% \\
    方案五 LoRA $r=16$ & 83.6\% & 16.4\% & 30.4\% & +5.6\% & +6.0\% & +3.8\% \\
    方案五 多候选重排序 & \textbf{97.8\%} & \textbf{2.2\%} & 37.4\% & +8.4\% & +11.1\% & +10.6\% \\
    \bottomrule
  \end{tabular}
\end{table}

由表~\ref{tab:evlm-results} 可以归纳出以下几点。首先，方案五（Qwen2-VL）显著优于方案一（纯文本 Qwen2.5 + LoRA）。Qwen2-VL 本身即为视觉语言模型，在预训练中已具备视觉-语言对齐能力；而普通 Qwen2.5 缺乏视觉理解基础，必须从头学习如何解读视觉前缀，因此输出稳定性较差。其次，LoRA 秩 $r=8$ 在 caption 类别命中率上优于 $r=16$，推测原因是 $r=8$ 在模型适配能力与过拟合风险之间取得了更合理的折中，而 $r=16$ 在小样本条件下可能导致生成稳定性下降。第三，多候选重排序策略在不改变模型参数的前提下，将有效输出率从 91.9\% 提升至 97.8\%，同时维持了较高的类别命中率。第四，方案五的所有变体在 real EEG 条件下的 caption 类别命中率均高于 vision-only 及 EEG 对照组的相应结果。

在此需要说明，推理过程中并未使用测试集的真实类别标签。多候选重排序完全基于候选 caption 的有效性、文本重复程度、长度以及与 A2/VTF 预测的 top-k 语义概念之间的相似度进行排序；caption 类别命中率仅用于最终评估环节，不参与候选选择的决策过程。

生成式模型的类别命中率仍然显著低于 A2 的 top-1 分类准确率。原因在于 A2 在 CLIP 语义空间中直接针对类别预测进行优化，目标函数和评估指标高度一致；而生成式模型在分类任务之外还需要学习语言生成能力，因此分类类指标的数值自然偏低。生成式模型的根本价值在于能够输出可理解的自然语言描述，适用于展示场景。

\subsection{定性示例分析}

为直观说明 real EEG 在退化条件下对语义理解的辅助作用，表~\ref{tab:qualitative} 展示了四个具有代表性的定性示例。每个示例同时给出了 vision-only、real EEG、shuffled EEG（打乱脑电）和 random EEG 四种模式下的 caption 输出。

\begin{table}[htbp]
  \centering
  \caption{定性 caption 示例对比}
  \label{tab:qualitative}
  \begin{tabularx}{\linewidth}{l l X X}
    \toprule
    真实类别 & 退化 & Vision-only caption & Real EEG caption \\
    \midrule
    reflex camera & lowres16 & a hand holding a cell phone & a black and white photo of an old camera \\
    canoe & lowres16 & a plane is flying through the sky & a red and yellow canoe in the water \\
    parachute & mixed & a group of people sitting on steps & a group of parachutes flying in the sky \\
    bolete mushroom & mixed & a man walking down a street at night & a mushroom in the ground among mushrooms \\
    \bottomrule
  \end{tabularx}

  \vspace{4pt}
  \begin{tabularx}{\linewidth}{l l X X}
    \toprule
    真实类别 & 退化 & Shuffled EEG caption & Random EEG caption \\
    \midrule
    reflex camera & lowres16 & a close up of a coffee machine & a close up of a shirt on a hanger \\
    canoe & lowres16 & a boat with boats in the background & a man holding a gun in his hand \\
    parachute & mixed & a group of people sitting around a fire & a woman ironing clothes \\
    bolete mushroom & mixed & a person holding flowers in the air & a girl ironing clothes on a board \\
    \bottomrule
  \end{tabularx}
\end{table}

从表~\ref{tab:qualitative} 中可以看出，在低分辨率和复合退化场景下，vision-only、shuffled EEG 和 random EEG 生成的 caption 与图像真实内容之间存在显著偏差（例如 canoe 被描述为 plane 或 gun，mushroom 被描述为 clothes 或 flowers）。相比之下，real EEG 生成的 caption 更接近于正确的物体类别，并且能够提供相对合理的自然语言描述。这一对比直观地印证了真实配对 EEG 在图像退化条件下对视觉语义理解具有辅助增强效果。

\subsection{实验结果讨论}

综合全部实验结果，可以总结出以下几点认识：

\textbf{A2 定量表现最佳的原因}：A2 的优化目标直接在 CLIP 语义空间中对齐类别预测，训练目标与评估指标之间高度一致；相比之下，生成式模型除了分类目标外还需兼顾语言生成的学习，在小样本条件下两类目标难以同时优化。因此 A2 的分类准确率大幅领先于所有生成式模型的 caption 类别命中率。

\textbf{VTF 定量表现弱于 A2 的原因}：token 级融合的设计初衷是为了满足生成式 VLM 的输入格式要求，而非追求最优的分类性能。在小样本条件下，token 级别的训练比 embedding 级别更难收敛，同时 VTF 仅采用单层 cross-attention，其交互深度不及 A2 的门控残差融合机制。因此 VTF 在分类指标上不及 A2 属于设计预期的正常结果。

\textbf{Qwen2-VL 方案优于 Qwen2.5 方案的原因}：Qwen2-VL 在预训练过程中已经完成了视觉-语言的对齐训练，其视觉 adapter 与语言 decoder 之间形成了有效的跨模态映射关系。纯文本 Qwen2.5 不具备视觉理解能力，在面对 EEG 增强的视觉前缀时，需要从零开始学习如何解读视觉信息，因此生成质量不够稳定。

\textbf{LoRA 秩的选择权衡}：$r=8$ 在模型适配能力与过拟合风险之间实现了更好的平衡；$r=16$ 引入了更多可训练参数，在小样本条件下可能导致生成稳定性降低，具体表现为有效输出率的下降。

\textbf{训练目标的选取策略}：以类别名称作为训练目标有助于提升 caption 类别命中率；以类别名称结合过滤后的 BLIP 伪 caption 作为训练目标则更有利于改善文本自然度和有效输出率。两类目标各有优势，本设计分别报告了各自的效果。

\textbf{实验的局限性}：所有结果均基于 Thought2Text 数据集的 40 个类别，尚不能直接推广到其他 EEG 数据集或更大规模类别体系。A2 的多种子评估仅采用了 3 个随机种子。生成式模型的类别评估依赖于 CLIP 文本原型的零样本分类能力，而该能力受到 CLIP 在具体类别上判别性能的限制。

\subsection{实验小结}

在本数据集及当前实验条件下可以得出以下结论。第一，A2 时间-频谱-空间语义融合模型在全部 18 组 seed-condition 组合中，real EEG 的准确率均高于 vision-only 以及 shuffled/random EEG，是本课程设计中最稳健的定量模型。第二，EEG 带来的性能增益与视觉退化程度之间存在正向关联。第三，基于 Qwen2-VL 的生成方案在强退化条件下实现了 97.8\% 的有效输出率，且 real EEG 相对于 vision-only 在 caption 类别命中率上提升了约 8 到 9 个百分点，但生成式模型在分类类指标上仍低于 A2 分类器。第四，控制实验表明，性能提升并非源自额外的输入通道或模型容量增加，而是依赖于 EEG 与图像之间的真实配对对应关系。
